\section{Introduction}
\label{sec:introduction}

Low Earth orbit (LEO) satellites extend Internet-of-Things (IoT)
connectivity to areas beyond the coverage of terrestrial networks.
The sporadic transmission of short packets by a large population of
terminals makes grant-free random access (GFRA) attractive for satellite
IoT, since active terminals can transmit without first obtaining an
uplink grant \cite{shen2025}. With limited pilot resources, the receiver
must identify active terminals and estimate their channels from
superimposed nonorthogonal pilot signals. This gives rise to the joint
activity detection and channel estimation (JADCE) problem
\cite{liu2018}. In LEO uplinks, the satellite motion and the different
propagation distances across its coverage area introduce user-dependent
delays and Doppler shifts, which further complicate pilot reception.

The Release~17 New Radio non-terrestrial network (NR-NTN) framework
assumes terminals with global navigation satellite system (GNSS)
capabilities. A terminal uses its GNSS position, satellite ephemeris,
and common timing advance information to perform uplink timing and
frequency pre-compensation \cite[Secs.~16.14.1 and 16.14.2.2]{3gpp38300}.
For low-cost, battery-powered IoT terminals, GNSS acquisition and the
associated compensation procedures impose additional hardware and
energy requirements \cite{shen2025}. This motivates GFRA schemes that
operate without terminal GNSS assistance. In the setting considered
here, a common reference delay is compensated, while the satellite
receiver handles the remaining differential delays and Doppler shifts.
These offsets can exceed a symbol duration and a subcarrier spacing,
respectively, and therefore need to be retained in the signal model
\cite{shen2025}.

Waveforms designed for doubly dispersive channels have been investigated
for satellite random access. Shen et al.\ proposed a spreading-based
multi-frame orthogonal time-frequency space (OTFS) transmission scheme
and a receiver for joint device identification, channel estimation, and
symbol detection in the presence of large differential delays and
Doppler shifts \cite{shen2025}. Affine frequency division multiplexing
(AFDM) provides another approach through the discrete affine Fourier
transform (DAFT), whose chirp parameters can be selected according to
the delay and Doppler spreads \cite{bemani2022}. For AFDM-based LEO
GFRA, Ding et al.\ developed a graph neural network that refines sparse
priors by learning correlations among delay--Doppler grid nodes
for JADCE \cite{ding2026}.

The choice of channel representation determines which dependencies a
receiver can exploit. A discrete delay--Doppler representation describes
the channel through path coefficients on a prescribed grid. Fractional
delays and Doppler shifts spread the contribution of a physical path
over multiple grid points, making the assumed support less accurate
\cite{shen2025,huang2026}. An alternative is to estimate the AFDM
effective channel directly \cite{huang2026}. In a multiantenna system,
each terminal's effective channel comprises a collection of matrices
that share path delays, Doppler shifts, and array responses. These
physical parameters induce dependencies among the matrix entries and
across receive antennas. Learning their joint distribution offers a way
to use these dependencies in channel estimation without prescribing a
discrete path support.

Score-based generative models learn a channel prior through the gradient
of the log density of Gaussian-perturbed channel samples. Arvinte and
Tamir applied such models to multiple-input multiple-output (MIMO)
channel estimation using posterior sampling \cite{arvinte2023}.
For LEO AFDM links, Huang et al.\ learned the score of the effective
channel and performed inference using likelihood-annealed Langevin
dynamics \cite{huang2026}. To incorporate generative priors into JADCE,
Cai et al.\ combined channel denoising with message passing (MP) in a
MIMO orthogonal frequency division multiplexing (OFDM) system
\cite{cai2026}. Their receiver uses first- and second-order score
networks to obtain channel moment estimates at each iteration, avoiding
an online reverse diffusion sampling procedure.

The Gaussian denoising model used within MP is an approximation to the
queries generated by a finite receiver. Wei and Yuan proposed learned
score-based turbo message passing (L-STMP) for single-user channel
recovery from spatially compressed measurements \cite{wei2026}.
Their method unfolds the receiver and adapts pretrained score networks
and calibration parameters through receiver training. For grant-free
access, this construction must additionally account for unknown
activity: the object entering the observation model is the product of
the activity indicator and the channel, rather than the active channel
alone. Consequently, channel reconstruction, activity probabilities,
and inter-layer feedback must be specified together.

In the GNSS-independent AFDM uplink, user-dependent phase increments
between pilot frames also enter the sensing matrix. The resulting
operator need not have orthonormal rows, even when the original pilots
are selected from a DFT matrix. We therefore use a linear update for
the general observation operator and combine conditional channel
denoising with a learned activity probability. The effective channel
estimate is their product. This construction retains the activity
mixture structure of \cite{cai2026} and the receiver adaptation approach
of \cite{wei2026}, while specifying the updates and supervision for
multiuser AFDM observations.

In this paper, we investigate generative-model-aided GFRA for a
GNSS-independent multiuser MIMO-AFDM satellite IoT system, where
single-antenna terminals communicate with a LEO satellite equipped with
a uniform planar array (UPA). Our main contributions are as follows.
\begin{enumerate}[label=\arabic*)]
\item We formulate a multi-frame AFDM uplink model with continuous path
delays and Doppler shifts. A chirp-periodic prefix (CPP) accommodates
the differential propagation delay, and the model separates
inter-frame phase evolution from the reference-frame effective
channel. Stacking the received pilots yields a multiuser observation
model for joint activity detection and full multiantenna channel
estimation.

\item We develop an unfolded MP receiver for JADCE under the general
AFDM observation operator. Each layer combines linear estimation with
a conditional channel denoiser and an activity classifier, and forms
the effective channel estimate by soft activity weighting. A calibrated
residual scale and Gaussian message damping specify the feedback to
the next layer without a separate second-order score network.

\item We train the receiver using active-user conditional channel
error and activity binary cross-entropy. A conditional-risk
decomposition explains the targets of these two losses and their
relation to the effective channel mean. Independently parameterized
layers are initialized from the same pretrained channel score network
and trained through layer adaptation and joint fine-tuning. We specify
the gradient paths, model-selection rule, and computational complexity.
\end{enumerate}

The conditional receiver accepts either supplied phase and noise
parameters or estimates obtained at the satellite. These settings are
distinguished in evaluation: training with supplied parameters isolates
the conditional receiver and does not establish the performance of the
complete receiver with estimated parameters.
% TWC numerical conclusions await independent evaluation of the new receiver.

The remainder of this paper is organized as follows.
Section~\ref{sec:system} presents the system model and problem
formulation. Section~\ref{sec:algorithm} develops the receiver, and
Section~\ref{sec:learning} describes the learning objectives.
Section~\ref{sec:overall} summarizes the algorithm and analyzes its
complexity. Numerical results and conclusions are presented in
Sections~VI and VII, respectively.
